\section{Related Work}

\subsection{Deep Learning in Dental Radiographs}
Deep learning has been widely applied to the automated interpretation of dental radiographs, including tooth detection and numbering and lesion detection and segmentation. Reviews indicate that convolutional neural networks can approach or even exceed expert performance in detection tasks on periapical and panoramic radiographs~\cite{Aminoshariae2021,Fontenele2025}. For tooth segmentation, Nguyen et al.~\cite{Nguyen2025} combined YOLO11 with a lightweight SAM to number and segment teeth in intraoral images, achieving a Dice coefficient of up to 0.983. For geometric measurement on periapical radiographs, Mao et al.~\cite{Mao2026} used YOLO to localise implants and graded peri-implant bone loss with image processing, and Lin et al.~\cite{Lin2026} combined a segmentation model with geometric algorithms to evaluate implant paths in edentulous regions with a mean angular error of only 2.79$^\circ$. These studies show that single-stage detection and segmentation models are clinically feasible for automated geometric measurement on periapical radiographs; however, they mainly address detection, classification, or angular measurement rather than millimetre-level length estimation when the pixel scale is unavailable.

\subsection{Working-Length and Tooth-Length Estimation}
For working-length estimation, Saghiri et al.~\cite{Saghiri2012a,Saghiri2012b} first used an artificial neural network with periapical radiographic features to determine the position of the file tip relative to the apical foramen, reaching accuracies of 93\% and 96\% in dry-skull and cadaver models, respectively; the latter exceeded the 76\% achieved by endodontists. More recently, Hsu et al.~\cite{Hsu2026} combined deep-learning segmentation with length conversion to estimate tooth length on periapical radiographs, which is the work most closely related to the present study.

\subsection{Image Enhancement and Small-Sample Regression}
The contrast of periapical radiographs is affected by exposure conditions and by the overlap of soft tissue and alveolar bone. Contrast-limited adaptive histogram equalisation (CLAHE) is a common preprocessing method for dental images that improves local contrast but may also amplify noise~\tbd{cite}. The morphological top-hat transformation estimates the slowly varying background through opening and closing and extracts bright and dark details smaller than the structuring element; it has been used to suppress non-uniform backgrounds and enhance fine structures~\tbd{cite}. For regression, neural networks are sensitive to random initialisation when samples are few, whereas GPR is a non-parametric Bayesian method whose hyperparameters can be estimated from the marginal likelihood and which provides predictive uncertainty, making it suitable for medical measurement problems with limited samples~\tbd{cite}.

\subsection{Research Gap}
Taken together, previous studies have demonstrated the feasibility of deep learning for detection, segmentation, and geometric measurement in dental radiographs, but their use for preoperative length estimation remains limited. Length-related studies either classify the position of the instrument rather than output a millimetre length~\cite{Saghiri2012a} or rely on a known pixel-to-millimetre scale and physical reference objects. For retrospectively collected data in which only de-identified exported images are available, the pixel scale is often missing and varies between images, preventing direct millimetre-level regression. In addition, the effects of pre-segmentation image enhancement and of the choice of regression model under small samples have rarely been compared systematically. The present study therefore extracts single teeth with a detector, segments them after top-hat enhancement, removes export-scale differences by long-side normalisation, and learns the mapping from normalised length to filling length with GPR, while comparing preprocessing methods and regression models one factor at a time.
